\documentclass[11pt]{article}
\usepackage[T1]{fontenc}
\usepackage{textcomp}
\usepackage[margin=0.75in]{geometry}
\usepackage{amsmath,amssymb,amsthm}
\usepackage{array,booktabs}
\usepackage{hyperref}
\setlength{\parindent}{0pt}
\newtheorem{theorem}{Theorem}
\theoremstyle{definition}
\newtheorem{definition}{Definition}
\title{Flow Proof Artifact: Nat-plus-cancel-right}
\author{Generated by \texttt{flow doc proof}}
\date{Flow Proof Book}
\begin{document}
\maketitle
Equal right summands can be cancelled when they sit on the right.\\[0.5em]
\textbf{Source.} peano — https://en.wikipedia.org/wiki/Cancellation\_property\\[0.5em]
\bigskip
\noindent{\large \textbf{Derived fact 1.} \textit{From b + a = c + a we deduce b = c} \hfill the natural numbers \cdot addition \cdot right cancellation holds}\par
\smallskip\noindent\textit{Coordinate: the natural numbers \cdot addition \cdot right cancellation holds}\par
\smallskip\noindent\textit{Source: peano/induction — Gries \& Schneider, Ch. 3}\par
\smallskip\noindent\textit{Built on: you can swap the order when you add, left cancellation holds, for addition on the natural numbers}\par
\medskip
\noindent\textbf{Goal.} From b + a = c + a we deduce b = c\par
\noindent$\displaystyle \forall a \in \mathbb{N} \forall b \in \mathbb{N} \forall c \in \mathbb{N}\quad b = c$\par
\medskip
\noindent
\begin{tabular}{@{} >{\raggedright\arraybackslash}p{0.06\textwidth} >{\raggedright\arraybackslash}p{0.47\textwidth} >{\raggedleft\arraybackslash}p{0.41\textwidth} @{}}
\textbf{\#} & \textbf{Proof} & \textbf{Mathematics} \\
\midrule
\textbf{1.} & We prove that right cancellation holds for addition on the natural numbers. & \\[0.45em]
\textbf{2.} & We split into exhaustive cases — the claim must hold in each one. & \\[0.45em]
\textbf{3.} & Case 1 (see step 2): suppose b + a  equals  c + a. & \\[0.45em]
\textbf{4.} & We invoke the derived fact governing addition on the natural numbers: you can swap the order when you add (instantiated for b, a). & $\displaystyle b + a = a + b$ \\[0.45em]
\textbf{5.} & We invoke the derived fact governing addition on the natural numbers: you can swap the order when you add (instantiated for c, a). & $\displaystyle c + a = a + c$ \\[0.45em]
\textbf{6.} & We invoke the derived fact governing addition on the natural numbers: left cancellation holds, for addition on the natural numbers (instantiated for a, b, c). & \\[0.45em]
\textbf{7.} & From step 3, step 4, step 5, and step 6, this implies b equals c. Together with the other cases (step 3), the goal is discharged. Hence proven. & $\displaystyle b = c$ \\[0.45em]
\bottomrule
\end{tabular}
\label{thm:Nat-addition-right_cancellation_holds}
\par\medskip\hrule
\end{document}
